\documentclass[11pt,a5paper]{book}
\usepackage[utf8]{inputenc}
\usepackage[T1]{fontenc}
\usepackage{geometry}
\geometry{a5paper, margin=1.2in, top=1.3in}
\usepackage{graphicx}
\usepackage{epigraph}
\usepackage{lettrine}
\usepackage{xcolor}
\usepackage{titlesec}
\usepackage{booktabs}
\usepackage{array}
\usepackage{enumitem}
\usepackage{quoting}
\usepackage{hyperref}
\usepackage{microtype}

% Fonts compatible with pdfLaTeX
\usepackage{libertine}
\usepackage[libertine]{newtxmath}
\usepackage{biolinum}
\renewcommand{\ttdefault}{cmtt}

% Color definitions
\definecolor{ashgray}{RGB}{178,178,178}
\definecolor{deepred}{RGB}{139,0,0}
\definecolor{fadedink}{RGB}{60,60,60}

% Author attribution commands
\newcommand{\defineauthor}[2]{%
  \expandafter\newcommand\csname #1\endcsname{\textsc{#2}}%
}
\defineauthor{gray}{The Gray Wanderer}
\defineauthor{braedon}{Duke Braedon Fenwood}
\defineauthor{baldvyn}{King Baldvyn Everblood}
\defineauthor{soren}{Søren the One-Eyed}
\defineauthor{valerius}{Lord Valerius Valvano}
\defineauthor{veralyn}{Queen Veralyn Everblood}

% Game mechanics shortcuts
\newcommand{\sb}[1]{\textbf{SB #1}}
\newcommand{\dv}[1]{\textbf{DV #1}}
\newcommand{\clock}[1]{\textbf{[#1]}}
\newcommand{\trait}[1]{\textbf{#1}}
\newcommand{\condition}[1]{\textit{#1}}

% Section formatting
\titleformat{\chapter}[display]
  {\huge\scshape\bfseries\centering}
  {\Large\chaptertitlename\ \thechapter}
  {20pt}
  {}
\titleformat{\section}
  {\Large\bfseries\scshape}
  {\thesection}{1em}{}
\titleformat{\subsection}
  {\normalsize\bfseries}
  {\thesubsection}{1em}{}

% Epigraph settings
\setlength{\epigraphwidth}{0.8\textwidth}
\renewcommand{\textflush}{flushright}

% Custom environment for "Witch's Closing"
\newenvironment{witchclosing}
  {\begin{quoting}[font={small\itshape}, leftmargin=2em, rightmargin=2em]\noindent\textsc{The Witch's Closing.}\par}
  {\end{quoting}}

% Custom environment for boxed mechanics
\newenvironment{mechanicbox}[1]
  {\begin{center}\fbox{\parbox{0.9\textwidth}{\textbf{#1}\par\smallskip}}}
  {\end{parbox}\end{center}}

\begin{document}

\frontmatter
\thispagestyle{empty}
\begin{center}
\vspace*{2cm}
{\Huge\scshape The Old Duke's Wars}\\[1em]
{\Large\textit{Vignettes from the Life of \braedon}}\\[2em]
{\large\gray}\\[0.5em]
{\normalsize\textit{Set down from the deathbed ledger}}\\[2em]
\vfill
{\small Tarlington, winter of 338 AR}\\[1em]
{\small Under the lamp, bread and salt.\\The ink is blood and lamp-black.}
\end{center}

\clearpage

\tableofcontents

\clearpage

\chapter*{A Note on This Writing}
\addcontentsline{toc}{chapter}{A Note on This Writing}

\lettrine{T}{his} ledger was found pressed between the leaves of a Canray slate, on the table beside Duke Braedon Fenwood's deathbed. The ink is his own blood---the last move of a game he could not finish. The pages are foxed, some stained by wine, others by tears. A single braid of Ykrul horsehair marks the place where the old man stopped writing.

\gray{} has added marginal notes in a drier hand. Where the duke's memory falters, we have bracketed the uncertainty. Where the truth is too cruel for ink, we have left a space.

This is not a chronicle of triumph. It is a record of what a house must bury to survive. And it ends where the next ledger begins---on the river Belworth, with blood in the water and a crown still warm.

\begin{mechanicbox}{The Deathbed Ledger}
This document serves as both narrative and \textbf{campaign frame}. The vignettes recount \braedon's life, but they also establish \textbf{Strings}, \textbf{Obligations}, and \textbf{Assets} for characters in the Fenwood bloodline.

\textbf{Required:} \textit{The Exile Road}, \textit{The Reckoning Bridge}, and the \textbf{Kon'reh} rules.
\end{mechanicbox}

\vspace{1em}
\noindent\textit{--- The Archive of the True Masons, Tarlington chapter, winter of 338 AR}

\mainmatter

\chapter{Frame: The Solar at Tarlington}

\epigraph{The fire burns low. Outside, the Belworth churns with winter melt. \braedon{} lies propped on pillows. A Canray slate rests on his chest. The pieces are not set.}{\gray}

\lettrine{T}{he} old duke is dying. His left arm is dead---eight years cold since the Sack of Lence---and his breath comes in shallow tides. But his eyes are clear, and his voice, when he speaks, carries the weight of sixty-six years of ledger-keeping.

You have been summoned to his solar. Not as servants. As witnesses. As the ones who must remember when he is ash.

\vspace{1em}
\noindent\textit{``Sit,''} he says. \textit{``I will not rise again. But I will speak. The road does not forgive a debt. It only reschedules it. You are the ledger now.''}

\vspace{1em}
\noindent He tells you seven stories. They are not in order. They are not kind. They are his confession, his instruction, and his final game of Canray---played with blood instead of wine.

\begin{mechanicbox}{Frame Mechanics: The Seven Vignettes}
Each vignette represents a \textbf{memory} that \braedon{} bequeaths. After each story, the players may ask one question. If they succeed on a \textbf{Wits + Insight} roll (\dv{3}), they gain a \textbf{String} on the subject of the vignette (e.g., \textit{``The Oath of the Yloka''} or \textit{``The Spared Thornhill''}).

When \braedon{} dies (at the end of the seventh vignette), any unspent Strings become \textbf{Obligations} owed by his heirs.
\end{mechanicbox}

\clearpage

\chapter{The Black Banner Years (Age Twenty--Two)}

\epigraph{I was a second son. My brother Conlin was the heir---clever, careful, born with a club foot that our father called a curse and I called an excuse.}{\braedon}

\lettrine{I}{} had no land, no title, no purpose. So I ran.

\baldvyn{} ran with me. He was also a second son---King Harald IV's second boy, not the heir. Harald the Bold was the elder; \baldvyn{} and I squired together under Oswin Everblood. We sharpened each other's blades and lied about each other's conquests. When I said, \textit{``I'm leaving,''} he said, \textit{``I'll bring the horses.''}

We joined the Black Banners---the ones who fought the Ykrul raiders in the steppe passes. Our captain was a woman called \soren{}, Linnic, scarred, and patient as frost. She taught us the first rule of war: \textit{``Count exits, not victims. The road remembers.''}

\section*{The First Blood}
Our first skirmish was a ford crossing. I was seventeen and terrified. \baldvyn{} was eighteen and pretending not to be. The Ykrul came out of the mist with wolf-tail standards. I killed a man---a boy, really---with a spear through the throat. He looked at me as if I had stolen something.

I did. I stole his tomorrows.

\baldvyn{} pulled me up. \textit{``Don't name him,''} he said. \textit{``Names are anchors. You'll carry him forever if you speak it.''}

I did not speak it. I still see his face.

\begin{witchclosing}
The first debt is always blood. It cannot be paid, only carried. The Black Banners do not keep ledgers of the dead; they keep braids of the living---to remind themselves who still owes the road a death.
\end{witchclosing}

\section*{The Bond}
Captain \soren{} matched us in the Kon'reh board. She played with stones, not pieces---river-smooth, each one a different shade of grey. \textit{``Play the player,''} she said, \textit{``not the position.''} \baldvyn{} beat me three times. I beat him twice. We never kept score. Friendship is not a ledger.

Later, in a camp by the Yloka, we swore an oath. Not on a book---on a braid of horsehair cut from the mane of a fallen Ykrul mare.

\textit{``If one of us becomes king, the other will be his hand. If one of us dies, the other will avenge him. If we both die, the road will remember.''}

He became king. I became his hand. I also married his sister, Aobrey. The road remembers.

\begin{mechanicbox}{Asset: The Oath of the Yloka}
If the party includes a Fenwood heir, they may invoke this oath once per campaign. When a Everblood is in danger, they may declare: \textit{``The Yloka remembers.''} Gain \sb{+2 dice} to any roll to save or avenge an Everblood, but mark \condition{Obligation 2} to the Ykrul Grey Ash clan (the braid's origin).
\end{mechanicbox}

\section*{The Return}
My brother Conlin died of a fever. I inherited the duchy. I came home to Tarlington with a Ykrul braid in my pocket, a Kon'reh board in my saddlebag, and Aobrey on a pillion behind me. My father met us at the gate.

\textit{``You smell of the steppe,''} he said. Then he kissed my bride's hand. \textit{``Good. Stinks of honest work.''}

\vspace{1em}
\noindent\textit{Thus endeth the first vignette. The fire cracks. \braedon{} drinks watered wine. His hand does not tremble yet.}

\clearpage

\chapter{The Thornhill Rebellion (Age Thirty--Six)}

\epigraph{The rape of Alayse Everblood---\baldvyn's youngest sister---was not a secret. Branton Thornhill made it a boast.}{\braedon}

\lettrine{T}{he} rape of Alayse Everblood was not a secret. Branton Thornhill made it a boast. He sent the king a lock of her hair wrapped in a marriage contract. \textit{``She came willingly,''} he wrote. \textit{``Your sister is a whore, and I am her lord.''}

King Harald IV rode to Brigh. He never came back. They killed him in a ditch, his own sword in his belly.

Then Harald the Bold---the elder brother, the heir---stormed the same walls. An arrow took his eye eleven days later. He lingered. He did not recover.

They crowned \baldvyn{} in a field, with a borrowed coronet and a dead brother's blood on his hands.

\section*{The Siege of Brigh}
\baldvyn{} summoned his banners. I brought eight hundred Fenwood spears---every man who could hold a pike, every woman who could ride. We encamped in the shadow of Branton's keep. The walls were high. The garrison was hungry. Branton mocked us from the parapet.

\textit{``Fenwood,''} he called down. \textit{``You were a mercenary. How much did they pay you to betray your people?''}

I did not answer. \baldvyn{} looked at me. \textit{``He wants rage. Give him geometry.''}

I spent the night tracing the keep's drainage sluice. By dawn, we had a tunnel. By noon, we were inside.

\begin{mechanicbox}{Tactical Lesson: The Geometry of Rage}
\braedon{} teaches: \textit{Rage is a weapon that wounds the wielder.} In combat, any PC who acts from \condition{Rage} gains \sb{+1 Harm} on their attacks but suffers \sb{--1 Position} on defense. The \textbf{Geometry} approach (Patience, Wits) treats the battle as a Kon'reh match: gain \sb{+1 die} to any roll that exploits terrain or supply lines instead of frontal assault.
\end{mechanicbox}

\section*{The Duel}
Branton met \baldvyn{} in the great hall. The fire was out. The tapestries had been cut down for bandages. They fought with longswords---no shields, no quarter. I held the door against Branton's guard.

\baldvyn{} was faster. Branton was stronger. It ended when \baldvyn{} caught Branton's blade on his gorget, shattered the steel, and drove his own sword through the lord's chest.

\textit{``For Alayse,''} he said. \textit{``For Harald. For my brother.''} He twisted the blade. Branton died with a smile.

\section*{The Child}
A servant woman fled the keep carrying a boy---Logan Thornhill, Branton's son, perhaps four years old. \baldvyn{} drew his dagger. I caught his wrist.

\textit{``He is a child.''}

\textit{``He is a seed.''}

\textit{``Then prune him gently, or the root will spread.''}

\baldvyn{} lowered the blade. He exiled the boy. The woman took him to Vhasia. Years later, that boy would kill Lyonel Vayle's father, and \veralyn{} would hunt him to a village called Coldmorrow, and she would burn it to ash.

\textit{``Mercy,''} \baldvyn{} said later, \textit{``is a debt. You just rescheduled it.''}

\vspace{1em}
\noindent\textit{\braedon{} sets down his cup. His hand trembles. The fire has burned low.}

\begin{witchclosing}
Mercy is the most expensive debt. It accrues interest not in gold, but in blood. The spared child grows to manhood, and the man remembers who held the blade that did not fall. This is why the Daughters of the Cord teach: \textit{Kill cleanly, or adopt fully. There is no middle ground.}
\end{witchclosing}

\clearpage

\chapter{The Longblood Defiance (Age Forty--Eight)}

\epigraph{Elianor of Adria was a warrior queen. She had fought beside \baldvyn, borne him children, ruled the marches when the king was away at war. But \baldvyn{} grew tired.}{\braedon}

\lettrine{E}{lianor} of Adria was a warrior queen. She had fought beside \baldvyn, borne him children, ruled the marches when the king was away at war. But \baldvyn{} grew tired. He took lovers---men and women both, openly, cruelly. Elianor's pride curdled into rebellion.

She raised the banners of Adria, Ethor, and the Bloods of Lence. Her son Valdais---or perhaps her brother? the lineage was tangled---rode at her side. \baldvyn{} called his host. I sat at the council table and watched two people I loved prepare to murder each other.

\section*{The Failed Mediation}
I rode to Elianor's camp under a flag of parley. She met me in a tent that smelled of wine and grief.

\textit{``You were his friend,''} she said. \textit{``You know what he is.''}

\textit{``I know what he was. He is still my king.''}

\textit{``Then you will die for him.''}

\textit{``I will die for the realm.''}

She laughed. It was not a kind sound. \textit{``There is no realm. There are only debts and ledgers.''}

I rode back. \baldvyn{} asked me what she said. I told him. He nodded. \textit{``Then let the ledgers speak.''}

\section*{The Battle of Charau}
I commanded the left flank. The ground was wet with autumn rain. The Bloods of Lence charged us three times. Three times we threw them back. Valdais---young, arrogant, his mother's fury in his eyes---broke through a gap in the line and rode straight for \baldvyn's standard.

I caught him with a mace. His horse went down. He crawled for his sword. I put my boot on his hand.

\textit{``Yield.''}

He spat at me.

I could have killed him. I should have killed him. But I saw Elianor's face in his. I pulled him up and marched him, stumbling, to \baldvyn.

\section*{The Pardon}
\baldvyn{} wanted execution. \textit{``He raised arms against his king.''}

\textit{``He is your son---or your half-brother. No one is certain. Does it matter?''}

\textit{``He will rebel again.''}

\textit{``Then let him owe his life to mercy. A debt unpaid is a leash.''}

\baldvyn{} chewed his rage for a long minute. Then he ordered Valdais confined to his estates. Elianor was exiled to a convent. The war ended.

Valdais never rebelled again. He became a quiet, loyal duke. And when \baldvyn{} died, Valdais wept.

\vspace{1em}
\noindent\textit{``Mercy,''} \braedon{} says, \textit{``is not weakness. It is a long loan.''}

\begin{mechanicbox}{The Long Loan}
\braedon{}'s mercy created a \textbf{Debt Clock} \clock{6} for Valdais. It never filled. When Valdais wept at \baldvyn's funeral, the clock was paid in full, and the debt was erased.

This is rare. Most debts of mercy fill their clocks and demand payment in the next generation.
\end{mechanicbox}

\clearpage

\chapter{The Sack of Lence (Age Sixty)}

\epigraph{The Vhasian king demanded homage. We had lands in Vhasia---old Valvano holdings, granted by the first Everbloods. We had never paid tribute.}{\braedon}

\lettrine{T}{he} Vhasian king demanded homage. We had lands in Vhasia---old Valvano holdings, granted by the first Everbloods. We had never paid tribute. He sent tax collectors. I sent them back with their papers shredded.

He declared war. I laughed. Then I raised the banners.

\section*{The Invasion}
My son Seniro commanded the vanguard. Wymund the Red---hot-headed, boisterous, secretly gentle. He rode with the Everbloods of Aberiel, with Harald and Elyas beside him. Young knights, hungry for glory. \veralyn{} was fifteen. She dressed as a page, hid in a supply wagon, presented herself at my tent.

\textit{``I will fight.''}

\textit{``Your father would kill me.''}

\textit{``My father is not here.''}

I assigned her to Seniro as a squire. She watched the streets of Lence run with blood.

\section*{The Siege}
We broke the walls on the third day. Street fighting---house to house, room to room. I lost my horse to a crossbow bolt. My left arm went dead. A Vhasian spearman knelt over me, raised his weapon.

Seniro cut him down. \textit{``Stay down, Father.''}

I stayed down. I watched from the mud as Seniro's men burned the palace.

\begin{mechanicbox}{Condition: The Dead Arm}
\braedon{} gains the permanent condition \condition{The Dead Arm} from this wound. All \textbf{Body} rolls involving two hands suffer \sb{--1 die}. However, he gains \sb{+1 die} to \textbf{Insight} rolls involving patience or suffering, having learned to wait in the mud.

This condition passes to any PC who inherits his mantle.
\end{mechanicbox}

\section*{The Last Thornhill}
Logan Thornhill---the boy I had spared---was a man now. He had fought for Vhasia, risen to captain, taken a new name. \veralyn{} tracked him to a village called Coldmorrow.

I did not see the killing. I saw the aftermath: smoke rising from a hundred thatched roofs, and \veralyn{} standing in the square with a mace dripping gore.

\textit{``He insulted my father,''} she said. \textit{``He killed my cousin.''}

\textit{``He was a child once.''}

\textit{``He was a monster grown.''}

I did not argue. I took her aside. \textit{``You will carry this weight forever. That is the price of the blade. Do not let it crush you. Let it make you sharp.''}

She wept. I held her. We did not speak of it again.

\vspace{1em}
\noindent\textit{\braedon{} touches his dead arm. ``The boy I spared killed the man I should have been.''}

\begin{witchclosing}
The spared enemy is a debt that always comes due. They do not forget the mercy, but they misremember it as weakness. When they return, they come not as grateful sons but as creditors, demanding payment in the coin of vengeance.
\end{witchclosing}

\clearpage

\chapter{The Dying King (Age Sixty--Four)}

\epigraph{Baldvyn called me to his bedside. He was wasted, yellowed, his breath sour with the canker that had eaten his stomach.}{\braedon}

\lettrine{B}{aldvyn} called me to his bedside. He was wasted, yellowed, his breath sour with the canker that had eaten his stomach. The physicians had bled him white. Nothing helped.

\textit{``I am dying,''} he said.

\textit{``You are old.''}

\textit{``I am sixty-four.''}

\textit{``So am I.''}

He laughed, then coughed blood. \textit{``Remember the Yloka? The ford?''}

\textit{``I remember.''}

\textit{``You saved my life.''}

\textit{``You saved mine first.''}

He pressed a letter into my hand. \textit{``Open it when I am gone. Not before.''}

\section*{The Tournament}
\baldvyn{} died a week later. His son Wyllmot---the Warrior, the Pious, the Brute---planned a grand tournament to honor the succession. \veralyn{} asked to compete. Wyllmot refused. She rode anyway.

She unhorsed three knights: Arin Tymber, my son Lerris, and Ouen Longblood. Ouen's horse died under him. He called her a whore. She dismounted and broke his jaw.

Tylan Vayle---quiet, patient, a wall of muscle---defeated her in the quarterfinal. He did not gloat. He offered her his hand. \textit{``You fought well, cousin.''}

She took it. She smiled for the first time in months.

\section*{The Letter}
I opened \baldvyn's letter after the tourney. It said only this:

\textit{``Watch over Wyllmot. He has my rage without my restraint. Watch over Veralyn. She has my mother's fire, and it will burn her if no one holds the bucket. You were my brother in all but blood. Do not let my blood destroy your name.''}

I burned the letter. The oath remains.

\vspace{1em}
\noindent\textit{``Wyllmot,''} \braedon{} says, \textit{``was a brute. But he was my king. I kept my oath.''}

\begin{mechanicbox}{The Unburned Oath}
Although \braedon{} burned the letter, the \textbf{Oath to Baldvyn} remains as a \condition{Persistent Debt}. Any Fenwood heir who fails to protect an Everblood from harm suffers \sb{--1 die} to all social rolls involving loyalty or honor until they perform a penance.

This debt cannot be paid off. It can only be passed to the next generation.
\end{mechanicbox}

\clearpage

\chapter{The Seven Week Uprising (Age Sixty--Five)}

\epigraph{Myrmis worshippers rose in the mountains. They killed lords, burned synods, declared a holy war against the Unity of Light.}{\braedon}

\lettrine{M}{yrmis} worshippers rose in the mountains. They killed lords, burned synods, declared a holy war against the Unity of Light. Wyllmot marched to storm their citadel---Eryas, a fortress carved into a cliff.

I counseled patience. \textit{``Let them starve. Winter is coming.''}

Wyllmot heard nothing. \textit{``They killed my cousin. They defiled a temple. I will have blood.''}

He led the assault himself. A rockslide caught the vanguard. Wyllmot died with a spear through his chest, still swinging his axe.

\section*{The Succession}
Baldvyn Brightstride---Wyllmot's son, young, untested---sat the throne. \veralyn{} rode for the mountains. I stayed in Tarlington with a garrison of old men and wounded.

\textit{``You are not a warrior anymore,''} my son Constano said.

\textit{``I am a duke. I am a strategist. I am a man who has buried too many friends.''}

\section*{The Conspiracy}
We learned later that Alaysha Renard---Wyllmot's widow, \veralyn's mother---had prayed to Myrmis on the eve of her wedding. She had sacrificed her favorite hound for three daughters and two sons. The goddess granted her wish---then demanded payment.

Alaysha had not paid. The goddess took Wyllmot.

I warned \veralyn: \textit{``Trust your mother not at all. She serves a god that demands sacrifice.''}

\veralyn{} said nothing. She had already known.

\vspace{1em}
\noindent\textit{``Some debts,''} \braedon{} says, \textit{``are not rescheduled. They are inherited.''}

\begin{witchclosing}
Gods are the ultimate creditors. They do not accept gold. They accept only what is precious---hounds, children, kings. Myrmis is patient. She will wait a generation for her payment. And when it comes due, it will be with interest.
\end{witchclosing}

\clearpage

\chapter{The Death of Brightstride (Two Weeks Ago – Age Sixty--Six)}

\epigraph{Baldvyn Brightstride was young, handsome, beloved. He had been married three months---to my granddaughter, Jaennis.}{\braedon}

\lettrine{B}{aldvyn} Brightstride was young, handsome, beloved. He had the charm of his father and the restraint of his grandmother. He had been married three months---to my granddaughter, Jaennis, a waif of a girl with chestnut hair and silent eyes.

He complained of a chill. He drank a posset. He died before dawn.

The physicians said fever. I said poison.

\section*{The Suspects}
The Renards had motive. The Longbloods had opportunity. The Vhasian Regency---still nursing the wounds of Lence---had means.

Jaennis wept at the funeral. I watched her. She did not weep for a husband. She wept for a future stolen.

\textit{``Did you kill him?''} I asked her later.

\textit{``I am a Fenwood. I do not answer that question.''}

\textit{``Answer it.''}

\textit{``He was kind. He was gentle. He did not deserve the crown.''}

\textit{``That is not an answer.''}

\textit{``It is the only one I have.''}

\section*{The Succession Crisis}
\veralyn{} is the named heir. But Valdais---Elianor's son, my spared captive---presses his claim. Aleric, Alaysha's nephew, whispers in dark corners. The realm teeters.

I have written \veralyn{} a letter. \textit{``Take the crown. Hold it. Let no one tell you that you are not worthy. You are a Fenwood. We are made of exile and iron.''}

I have not sent it. I will send it with you.

\vspace{1em}
\noindent\textit{\braedon{} hands a sealed parchment to the senior PC. ``Deliver this. Then come back. The river is rising. The next war begins at dawn.''}

\begin{mechanicbox}{The Succession Crisis Clocks}
As \braedon{} dies, the following clocks begin:

\textbf{Valdais's Claim \clock{6}}: Advances when his legitimacy is discussed. When full, he marches on the capital.

\textbf{Veralyn's Coronation \clock{4}}: Advances when her claim is supported. When full, she is crowned---but the realm is divided.

\textbf{The Poison Investigation \clock{8}}: Advances when clues are found. When full, the killer is revealed (Jaennis, Alaysha, or Vhasian agents---GM's choice).
\end{mechanicbox}

\clearpage

\chapter{Epilogue: The Road Reschedules}

\lettrine{T}{he} fire is ash. \braedon{} lies still. His chest rises once, twice---then stops.

A servant draws the sheet over his face. The Canray slate falls from his hand. The pieces scatter.

One---a Blue---comes to rest against the hearth. The rest are lost in shadow.

\vspace{1em}
\noindent \gray{} picks up the slate. The unfinished game is a puzzle. \textit{``He was playing for time. He always played for time.''}

\vspace{1em}
\noindent Outside, the Belworth churns. Riders are crossing. The road does not forgive a debt. It only reschedules.

\vspace{2em}
\noindent\textit{--- To be continued in \textbf{``Blood on the Belworth.''}}

\clearpage

\appendix
\chapter*{Appendix: The Canray Board \& Mechanical Integration}
\addcontentsline{toc}{chapter}{Appendix: The Canray Board}

\section*{The Unfinished Game}

\braedon's Canray slate represents the \textbf{Succession Crisis}. The GM should use actual Kon'reh pieces or a diagram:

\begin{center}
\begin{tabular}{@{}lp{8cm}@{}}
\toprule
\textbf{Piece} & \textbf{Significance} \\
\midrule
Blue (Throne) & Isolated, with few supporting pieces---\veralyn's precarious position. \\
Oranges (Allies) & Two forks: one aligned with Blue, one distant---the divided Viterran/Vhasian courts. \\
Reds (Fenwood) & Defensive wall along the home rank, with a single gap (the Belworth crossing). \\
Green (Runner) & Locked behind enemy lines---the message \braedon{} wants delivered. \\
\bottomrule
\end{tabular}
\end{center}

\textbf{The Solution:} Sacrifice the Blue to free the Green. A metaphor for \braedon{} dying so \veralyn{} can escape the board (the throne) and live.

\section*{Inheritable Conditions}

Any PC who is a Fenwood heir may inherit:

\begin{itemize}[leftmargin=*]
\item \condition{The Dead Arm}: --1 die to two-handed Body rolls; +1 die to patience-based Insight rolls.
\item \condition{The Oath of the Yloka}: Once per campaign, +2 dice to save an Everblood; Obligation 2 to Grey Ash.
\item \condition{The Spared Thornhill}: Logan Thornhill's descendants (if any survived Coldmorrow) have a String on you.
\end{itemize}

\clearpage

\chapter*{Colophon}
\addcontentsline{toc}{chapter}{Colophon}

\lettrine{T}{his} book was written in the solar at Tarlington, on paper that had once been a grain manifest. The ink is lamp-black cut with the ash of a burned confession. The binding is horsehair, tied with nine knots.

The ninth knot is not cut.

The old duke is dead. The road is still south. The river rises.

\vspace{1em}
\noindent\textit{--- \gray, at the edge of the Belworth, dawn of the Crisis.}

\end{document}

